\documentclass[a4paper]{article}
% Español (México) con XeLaTeX: fontspec para fuentes Unicode, polyglossia
% para la división silábica y los nombres del documento en la variante mexicana.
\usepackage{fontspec}
\usepackage{polyglossia}
\setdefaultlanguage[variant=mexican]{spanish}
% Los nombres chinos van como «pinyin (original)»: xeCJK da a los caracteres
% Han su propia fuente; sin ella XeLaTeX los omite con un aviso.
% FandolSong no cubre algunos caracteres de nombres (U+73FA); WenQuanYi Zen Hei sí.
\usepackage[AutoFallBack=true]{xeCJK}
\setCJKmainfont{FandolSong-Regular.otf}
\setCJKfallbackfamilyfont{\CJKrmdefault}{WenQuanYi Zen Hei}
% polyglossia no traduce los nombres de listings.
\AtBeginDocument{\providecommand{\lstlistingname}{}\renewcommand{\lstlistingname}{Listado}%
  \providecommand{\lstlistlistingname}{}\renewcommand{\lstlistlistingname}{Índice de listados}}
\usepackage{amsmath, amssymb}
\usepackage{graphicx}
\usepackage[margin=2.35cm]{geometry}
\usepackage[most]{tcolorbox}
\usepackage{booktabs}
\usepackage{tabularx}
\usepackage{longtable}
\usepackage{array}
\usepackage{float}
\usepackage{xcolor}
\usepackage{hyperref}
\usepackage{fancyhdr}
\setCJKmainfont[AutoFakeBold=2.4]{AR PL UMing CN}
\setCJKsansfont[AutoFakeBold=2.4]{Droid Sans Fallback}
\setCJKmonofont[AutoFakeBold=2.4]{Droid Sans Fallback}
\hypersetup{colorlinks=true, linkcolor=blue!55!black, urlcolor=blue!60!black}
\graphicspath{{./}{figures/}}
\setlength{\parskip}{0.55em}
\setlength{\parindent}{2em}
\setlength{\emergencystretch}{3em}
\renewcommand{\arraystretch}{1.18}
\newtcolorbox{knowledgebox}[1]{enhanced, colback=blue!5!white, colframe=blue!70!black, colbacktitle=blue!70!black, coltitle=white, fonttitle=\bfseries, title={#1}, sharp corners, breakable}
\newtcolorbox{importantbox}[1]{enhanced, colback=yellow!10!white, colframe=yellow!75!black, colbacktitle=yellow!75!black, coltitle=black, fonttitle=\bfseries, title={#1}, sharp corners, breakable}
\newtcolorbox{warningbox}[1]{enhanced, colback=red!5!white, colframe=red!70!black, colbacktitle=red!70!black, coltitle=white, fonttitle=\bfseries, title={#1}, sharp corners, breakable}
\pagestyle{fancy}
\fancyhf{}
\fancyfoot[C]{\thepage}
\fancyfoot[R]{\footnotesize\color{gray}\href{https://github.com/hqhq1025/ai-course-notes}{AI Course Notes}}
\renewcommand{\headrulewidth}{0pt}
\renewcommand{\footrulewidth}{0pt}
\newcommand{\videourl}{https://www.youtube.com/watch?v=JxEetUlV9RA}
\newcommand{\repourl}{https://github.com/hqhq1025/ai-course-notes}

\begin{document}
\begin{titlepage}
\centering
{\Large Notas didácticas de entrevista a fondo\par}
\vspace{1.0cm}
{\huge\bfseries Li Yifan (李一帆) relata 11 años de emprendimiento en lidar}\par
\vspace{0.5cm}
{\Large Zhang Xiaojun conversa con Li Yifan, cofundador y CEO de Hesai (禾赛)\par}
\vspace{0.35cm}
{\large Elaborado a partir del video público de Zhang Xiaojun Podcast, la transcripción hecha con GPU, la estructura de capítulos y diagramas conceptuales propios\par}
\vspace{0.3cm}
{\large 2025-08-07\par}
\vspace{0.85cm}
\includegraphics[width=0.88\textwidth,height=0.42\textheight,keepaspectratio]{cover.jpg}\par
\vfill
\begin{tcolorbox}[width=0.92\textwidth, colback=black!2!white, colframe=black!55, sharp corners]
\textbf{Título del video}: 111. Li Yifan relata 11 años de emprendimiento en lidar: si lo piensas con detenimiento, ¿de dónde vienen las oportunidades de la industria? Son oportunidades del país y de la nación\par
\textbf{Canal del video}: Zhang Xiaojun Podcast\par
\textbf{Duración del video}: 03:08:32\par
\textbf{Enlace del video}: \href{\videourl}{\nolinkurl{\videourl}}\par
\textbf{Notas sobre el material}: YouTube no tiene subtítulos disponibles; el texto se elaboró a partir de la transcripción con faster-whisper large-v3 en CUDA, los capítulos del video, la descripción del video, la portada y diagramas didácticos propios. Las tomas fijas de la entrevista no se repiten en el texto.\par
\textbf{Página del proyecto}: \href{\repourl}{\nolinkurl{\repourl}}\par
\end{tcolorbox}
\end{titlepage}

\tableofcontents
\newpage

\section{Introducción: por qué la historia de una empresa emergente de LiDAR pertenece a un curso sobre la industria tecnológica}

Esta sección explica primero por qué este episodio entra en la línea general de IA e internet. EP111 no es una entrevista con una empresa de modelos, sino un caso de emprendimiento de tecnología profunda en la cadena de suministro automotriz. En los últimos diez años, los vehículos de nueva energía en China surgieron prácticamente desde cero; las marcas de vehículos son muy visibles, pero detrás de ellas hay eslabones poco visibles: sensores, chips, cadena de suministro, manufactura de grado automotriz y relaciones con las armadoras. La historia del LiDAR de Hesai (禾赛) muestra cómo una empresa de tecnología profunda pasa de una tecnología de laboratorio a la producción en serie de grado automotriz.

Las palabras clave de este episodio no son «tráfico», sino «reducción de costos, pedidos, fijación de precios, deserción de clientes, armadoras, capital y oportunidad industrial nacional». El episodio responde a una pregunta: cuando la innovación tecnológica en China pasa de la innovación en modelos de negocio de internet a la innovación de frontera en tecnología profunda, ¿qué capacidades necesitan adquirir los emprendedores tecnológicos?

\begin{figure}[H]
\centering
\includegraphics[width=0.96\textwidth]{hesai-timeline.png}
\caption{Los 11 años de trayectoria emprendedora de Hesai: de la tecnología de laboratorio y el primer pedido grande hasta la entrada en el núcleo de la industria automotriz. Diagrama conceptual de elaboración propia, basado en los capítulos del video y en la conversación de 00:03:40--03:02:16.}
\end{figure}

\begin{knowledgebox}{Lectura de la figura: emprender en tecnología profunda no es un invento aislado}
Una empresa de LiDAR tiene que resolver al mismo tiempo la tecnología, los costos, la manufactura, la validación con clientes, el financiamiento, la fijación de precios y la ventana industrial. Si falla cualquiera de estos eslabones, es difícil que la tecnología se convierta en una posición dentro de la industria.
\end{knowledgebox}

\begin{importantbox}{Tesis central de este episodio}
Las oportunidades del emprendimiento en tecnología profunda suelen surgir de la suma de la cadena industrial y las capacidades nacionales: un avance tecnológico aislado no basta; tiene que coincidir con la cadena de suministro, las armadoras, el ciclo de capital, la capacidad de producción en serie y la ventana industrial.
\end{importantbox}

\begin{knowledgebox}{Emprendimiento en tecnología profunda frente a emprendimiento en internet}
\begin{tabularx}{\textwidth}{p{0.20\textwidth}p{0.34\textwidth}X}
\toprule
Dimensión & Emprendimiento en internet & Emprendimiento en tecnología profunda \\
\midrule
Activos centrales & Producto, tráfico, efectos de red & Tecnología, cadena de suministro, manufactura, certificación de clientes. \\
Velocidad de iteración & Lanzamiento rápido, pruebas A/B rápidas & Prototipos, pruebas, certificación; ciclos largos hasta la producción en serie. \\
Costo del fracaso & Se puede revertir, se puede rediseñar & Riesgos altos en materiales, inventario, calidad y clientes. \\
Uso del financiamiento & Crecimiento, investigación y desarrollo, ampliación del equipo & Investigación y desarrollo, equipo, cadena de suministro, entregas e inventario. \\
Puntos de inflexión clave & Crecimiento y retención de usuarios & Nominación por parte del cliente, producción en serie, curva de costos y recompra. \\
\bottomrule
\end{tabularx}
\end{knowledgebox}

\subsection{Resumen de la sección}

EP111 es una lección sobre el emprendimiento en tecnología profunda. Sitúa el LiDAR, más allá de ser un sensor, en el contexto general de la cadena industrial de los vehículos de nueva energía y del emprendimiento chino en tecnología profunda.
\section{Lidar: los ojos activos de autos y robots}

Esta sección explica primero qué es el lidar. Li Yifan (李一帆) lo compara con los ojos de los autos y los robots. La cámara también es un ojo, pero recibe la luz de forma pasiva; el lidar emite luz láser de forma activa y mide la distancia a partir del tiempo de vuelo de la luz, por lo que obtiene la información de distancia de manera más directa. En la conducción autónoma, la robótica y los escenarios relacionados con la seguridad, la estimación de distancias suele ser una cuestión de vida o muerte.

\begin{figure}[H]
\centering
\includegraphics[width=0.94\textwidth]{lidar-active-eye.png}
\caption{Lidar: un ojo activo. La cámara observa el mundo de forma pasiva; el lidar emite luz de forma activa y mide distancias. Diagrama conceptual propio, elaborado a partir de la conversación de 00:03:44--00:04:46.}
\end{figure}

\begin{knowledgebox}{Lectura de la figura: el sentido de la medición activa de distancias}
La cámara puede estimar distancias a partir de la imagen, pero el lidar mide directamente el tiempo de ida y vuelta de la luz y obtiene información de distancia más estable. Para un auto o un robot no basta con saber «qué se ve»; también hay que saber «a qué distancia está de mí».
\end{knowledgebox}

\subsection{Reducción de costos del 99.5\,\%: de instrumento científico a componente de grado automotriz}

La sección anterior explicó por qué el lidar es importante; esta aborda una cuestión más industrial: cómo pudo pasar de ser un costoso instrumento científico a un componente de grado automotriz. El juicio central es que solo cuando el costo se replantea por completo el lidar pasa de unos cuantos experimentos y proyectos de gama alta a la cadena de suministro principal de la industria automotriz. El título del video destaca la reducción del 99.5\,\% en el costo del lidar. Detrás de esa cifra no hay una simple rebaja en las compras, sino la acción conjunta de la arquitectura del producto, los componentes centrales, los chips, el desarrollo propio, la cadena de suministro y la manufactura a escala. Ahí está también lo más difícil de emprender en hardware: no solo hay que lograr el desempeño, sino llevarlo a un costo que permita la producción en serie, el mantenimiento y la aceptación de los fabricantes de automóviles.

\begin{figure}[H]
\centering
\includegraphics[width=0.96\textwidth]{cost-down-stack.png}
\caption{Ruta de reducción de costos del 99.5\,\%: pasar del prototipo de laboratorio a la producción en serie de grado automotriz exige que el diseño, los chips, la cadena de suministro y la manufactura bajen el precio en conjunto. Diagrama conceptual propio, elaborado a partir de la conversación de 00:03:40--00:12:05.}
\end{figure}

\begin{knowledgebox}{Lectura de la figura: la ruta de reducción de costos se lee de izquierda a derecha}
La etapa de prototipo resuelve «si se puede construir»; la reducción de costos por diseño resuelve «si se puede rediseñar en una forma barata»; el desarrollo propio de chips resuelve «si el costo y el desempeño centrales están bajo control»; la cadena de suministro resuelve «si se puede comprar y fabricar de forma estable»; solo en la etapa de producción en serie se entra de verdad en la industria automotriz.
\end{knowledgebox}

\begin{knowledgebox}{Términos clave: industrialización del lidar}
\begin{tabularx}{\textwidth}{p{0.22\textwidth}p{0.32\textwidth}X}
\toprule
Concepto & Significado & Por qué importa \\
\midrule
Medición activa de distancias & Emitir luz láser y medir el tiempo de retorno & Proporciona distancias precisas y la estructura espacial. \\
Grado automotriz & Cumplir los requisitos automotrices de confiabilidad, temperatura, vida útil y seguridad & Es el umbral para pasar del prototipo a la producción en serie. \\
Costo BOM & Costo de la lista de materiales & Determina si el hardware puede llegar a modelos de gran volumen. \\
Desarrollo propio de chips & Integrar y especializar los componentes centrales & Es un medio importante para reducir costos drásticamente y optimizar el desempeño. \\
\bottomrule
\end{tabularx}
\end{knowledgebox}

\begin{knowledgebox}{Reducir costos no es una acción única}
\begin{tabularx}{\textwidth}{p{0.20\textwidth}p{0.34\textwidth}X}
\toprule
Nivel de reducción & Acciones concretas & Riesgos \\
\midrule
Diseño de la solución & Cambiar la trayectoria óptica, la estructura, el método de escaneo y la arquitectura del sistema & Puede sacrificarse desempeño o confiabilidad. \\
Componentes centrales & Láser, receptor, escáner, integración en chip & Ciclos de desarrollo largos y alto riesgo de rendimiento de fabricación. \\
Procesos de manufactura & Ensamble, calibración y pruebas automatizados & El aumento gradual de la producción en serie saca a la luz muchos problemas. \\
Cadena de suministro & Compras a escala y sustitución por proveedores nacionales & Hay que mantener la calidad y la consistencia. \\
Adaptación a modelos & Diseño conjunto con la plataforma del fabricante de automóviles & Los cambios en los requisitos del cliente obligan a rediseñar una y otra vez. \\
\bottomrule
\end{tabularx}
\end{knowledgebox}

\begin{importantbox}{Por qué una reducción de costos del 99.5\,\% es un acontecimiento industrial}
Si un componente solo baja 10\,\% de precio, suele tratarse de una optimización de compras; si baja 99.5\,\%, normalmente significa que la forma del producto, los dispositivos centrales, el método de fabricación y el mercado de aplicación se han replanteado por completo. El paso del lidar de sensor costoso a producción en serie de grado automotriz es justamente ese replanteamiento.
\end{importantbox}

\begin{knowledgebox}{Implicaciones de mercado de la reducción de costos}
La reducción de costos hace que el lidar pase de «lo pueden usar unos cuantos vehículos de prueba de gama alta» a «lo pueden considerar más modelos de producción en serie». Esto cambia la estructura de clientes, la forma de vender, los requisitos de calidad y la lógica de competencia: venderle a un equipo de investigación y desarrollo y venderle a la plataforma de producción en serie de un fabricante de automóviles son dos negocios completamente distintos.
\end{knowledgebox}

\subsection{Resumen de la sección}

El lidar pasa de ser «una buena tecnología» a «una oportunidad industrial» gracias a que se cumplen a la vez el valor de la medición activa de distancias y la capacidad de reducir costos de forma extrema. Sin reducción de costos no hay producción en serie de grado automotriz; sin producción en serie de grado automotriz no hay lugar en la industria.
\section{Estructura de la empresa: tres socios con participaciones iguales, financiamiento y el primer pedido grande}

La sección anterior trató la tecnología; esta examina la organización y el capital. Un rasgo poco común de los inicios de Hesai (禾赛) es que los tres fundadores se repartieron las acciones en partes iguales. Para una empresa emergente, esto expresa una confianza sólida, pero también impone exigencias de gobierno corporativo. Una empresa emergente de tecnología profunda tiene ciclos largos, necesita mucho capital y enfrenta gran incertidumbre técnica y comercial; si los fundadores no se comprometen entre sí a largo plazo, difícilmente superarán las etapas bajas de los primeros años.

\begin{figure}[H]
\centering
\includegraphics[width=0.94\textwidth]{three-founder-equity.png}
\caption{Tres socios con participaciones iguales: detrás de esta estructura accionaria poco común están la confianza, la complementariedad de capacidades y el compromiso a largo plazo. Diagrama conceptual propio, elaborado a partir de la conversación entre 00:32:13 y 00:43:35.}
\end{figure}

\begin{knowledgebox}{Lectura de la figura: repartir las acciones en partes iguales no es romanticismo}
El reparto en tres partes iguales puede fortalecer el sentido de comunidad, pero exige que el equipo fundador tenga mecanismos sólidos de confianza mutua, comunicación y toma de decisiones. De lo contrario, al llegar a las etapas de financiamiento, cambio de rumbo y reparto de beneficios, los riesgos de gobierno corporativo se agravan.
\end{knowledgebox}

\subsection{El financiamiento no es un relato, sino un puente de flujo de efectivo}

Lo anterior trató la estructura accionaria; esta sección pasa al ritmo del capital y explica por qué a una empresa de tecnología profunda no le basta con presentar una visión. Para ella, financiarse no consiste solo en presentar una visión: el capital debe sostener la investigación y el desarrollo, los prototipos, la cadena de suministro, las entregas y la validación con clientes. Li Yifan (李一帆) habló del ritmo de financiamiento y de la presión sobre el flujo de efectivo, lo que muestra que una empresa emergente de tecnología profunda no puede depender solo del relato ante los inversionistas: al final tiene que superar la prueba de los pedidos y de la cobranza.

\begin{figure}[H]
\centering
\includegraphics[width=0.96\textwidth]{financing-cashflow.png}
\caption{Financiamiento y flujo de efectivo: una empresa emergente de hardware no puede limitarse a contar un relato; tiene que superar los prototipos, los pedidos y el flujo de efectivo. Diagrama conceptual propio, elaborado a partir de la conversación entre 00:43:35 y 00:49:02 y entre 01:10:06 y 01:20:47.}
\end{figure}

\begin{knowledgebox}{Lectura de la figura: el financiamiento solo compra tiempo dentro de la cadena de validación}
En la figura, cada paso, desde el relato tecnológico hasta la expansión de la producción en serie, se acerca más al mercado real. El financiamiento puede sostener la investigación y el desarrollo, pero solo los pedidos y la cobranza demuestran que el cliente está dispuesto a asumir un riesgo real; la producción en serie, a su vez, indica que la organización y la cadena de suministro empiezan a superar la prueba.
\end{knowledgebox}

\subsection{El primer pedido de 20 millones}

El financiamiento solo compra tiempo; los pedidos son los que demuestran que existe un mercado. Esta sección examina el significado del primer pedido grande, de 20 millones. Ese primer pedido de 20 millones validó al equipo en el mercado por primera vez y, al mismo tiempo, multiplicó la presión de entrega. Para una empresa de tecnología profunda, un pedido grande no es un simple ingreso, sino un hecho que construye credibilidad: que un cliente esté dispuesto a confiarte un escenario real significa que el producto empieza a salir del laboratorio.

\begin{figure}[H]
\centering
\includegraphics[width=0.96\textwidth]{first-big-order.png}
\caption{El primer pedido de 20 millones: un pedido grande valida por primera vez en el mercado al equipo técnico y, al mismo tiempo, multiplica la presión de entrega. Diagrama conceptual propio, elaborado a partir de la conversación entre 00:49:02 y 00:55:45.}
\end{figure}

\begin{warningbox}{Un pedido grande también multiplica los riesgos}
Si falla la entrega de un pedido grande de hardware, se dañan a la vez el flujo de efectivo, la confianza del cliente y la reputación en la industria. Firmar el contrato es solo el comienzo; la prueba real está en la entrega, la estabilidad y el servicio posventa.
\end{warningbox}

\begin{knowledgebox}{Cadena de validación temprana en tecnología profunda}
\begin{tabularx}{\textwidth}{p{0.20\textwidth}p{0.34\textwidth}X}
\toprule
Etapa & Qué se valida & Evidencia principal \\
\midrule
Prototipo & Si la tecnología funciona & Demostraciones, datos de pruebas, reconocimiento de expertos. \\
Pedido pequeño & Si el cliente está dispuesto a probarlo & Contrato piloto, retroalimentación en campo. \\
Pedido grande & Si el cliente está dispuesto a asumir riesgos & Presupuesto real, presión de entrega real. \\
Recompra & Si el producto es realmente útil & El cliente sigue comprando y amplía la escala. \\
Producción en serie & Si puede integrarse en la cadena industrial & Entrega estable, sistema de calidad, reducción de costos. \\
\bottomrule
\end{tabularx}
\end{knowledgebox}

\begin{knowledgebox}{Diferencias entre financiamiento, pedidos y cobranza}
\begin{tabularx}{\textwidth}{p{0.20\textwidth}p{0.34\textwidth}X}
\toprule
Concepto & Qué demuestra & Qué no demuestra \\
\midrule
Financiamiento & Que los inversionistas creen en las posibilidades futuras & No demuestra que el cliente vaya a comprar. \\
Pedido & Que el cliente está dispuesto a destinarte su presupuesto & No demuestra que puedas entregar de forma estable. \\
Cobranza & Que el cliente acepta la entrega y cumple con el pago & No demuestra que esté resuelto el costo a escala. \\
Recompra & Que el cliente está dispuesto a seguir comprando & Solo entonces se empieza a acercar al ajuste real entre producto y mercado. \\
\bottomrule
\end{tabularx}
\end{knowledgebox}

\subsection{Resumen de la sección}

La estructura organizativa y de capital de los inicios de Hesai muestra que una empresa emergente de tecnología profunda debe resolver al mismo tiempo la confianza dentro del equipo, el ritmo de financiamiento, los pedidos de los clientes y la credibilidad en las entregas. La tecnología es solo el boleto de entrada a la mesa de juego.
\section{De «se acabó» al cambio de proveedor de los clientes}

Esta sección entra en el punto más bajo de la empresa y en el punto de inflexión. En el video hay un segmento sobre «querer decir que se acabó»: el equipo tenía algo de tecnología láser, pero no un mercado lo bastante atractivo. Ese momento es decisivo, porque obliga al equipo a dejar de ser emprendedores tecnológicos y convertirse en verdaderos emprendedores de mercado: ya no se preguntan solo qué puede hacer la tecnología, sino por qué compra el cliente, cuándo compra y a qué precio.

\begin{importantbox}{El primer cambio de rumbo del emprendimiento en tecnología profunda}
Pasar de «tengo la tecnología» a «por qué el cliente tiene que comprar» es el cambio de rumbo decisivo del emprendimiento en tecnología profunda. Lo primero responde a una lógica de investigación y desarrollo; lo segundo, a una lógica de industria.
\end{importantbox}

\subsection{La lógica de precios: costo, valor y posición estratégica}

La subsección anterior planteó que el emprendimiento en tecnología profunda debe pasar de una lógica tecnológica a una lógica centrada en el cliente; esta subsección lleva ese cambio a la fijación de precios. El precio de un producto de hardware no se calcula como costo más margen. Debe considerar el costo de la BOM, el valor de seguridad para el cliente, el precio de los productos competidores que podrían sustituirlo, la reducción de costos esperada y la posición estratégica que da entrar en las armadoras. Si el precio inicial es incorrecto, se pueden perder clientes; si es demasiado bajo, quizá no alcance para sostener las entregas y la iteración.

\begin{figure}[H]
\centering
\includegraphics[width=0.96\textwidth]{pricing-logic.png}
\caption{La lógica de precios del hardware: el precio no es costo más margen, sino valor para el cliente, riesgo y posición en el mercado. Diagrama conceptual propio, elaborado a partir de la conversación de 01:20:47--01:38:15.}
\end{figure}

\begin{knowledgebox}{Lectura de la figura: el precio es un lenguaje estratégico}
El precio de un producto de hardware le comunica al cliente tres cosas a la vez: cuánto vale el producto, cuánto riesgo asume el proveedor y si en el futuro podrá producirse a escala. Si el precio inicial se calcula solo con base en el costo, se subestima el valor estratégico de entrar en la cadena de suministro de las armadoras.
\end{knowledgebox}

\begin{knowledgebox}{Las cinco variables del precio del hardware}
\begin{tabularx}{\textwidth}{p{0.20\textwidth}p{0.34\textwidth}X}
\toprule
Variable & Significado & Efecto \\
\midrule
Costo & Materiales, manufactura, rendimiento de producción, servicio posventa & Define el piso. \\
Valor & Mejora en seguridad, experiencia y desempeño & Define cuánto está dispuesto a pagar el cliente. \\
Competencia & Proveedores actuales y soluciones alternativas & Define el margen para entrar. \\
Escala & Reducción de costos con la producción en serie futura & Define la estrategia de precios a largo plazo. \\
Estrategia & Si se entra en la cadena de suministro de un cliente clave & Define si se puede sacrificar la utilidad a corto plazo. \\
\bottomrule
\end{tabularx}
\end{knowledgebox}

\begin{warningbox}{La trampa de los precios bajos en hardware}
Un precio bajo puede acelerar que los clientes cambien de proveedor, pero también puede dejar a la empresa con margen bruto negativo, carga de servicio posventa y presión sobre la calidad durante mucho tiempo. El precio del hardware debe servir a la vez para entrar al mercado y para mantener viva a la organización; no puede perseguir solo «ganar al cliente».
\end{warningbox}

\subsection{Los clientes empiezan a cambiar de proveedor}

El precio muestra cómo el proveedor expresa su valor; esta subsección examina cómo votan de verdad los clientes. Que un cliente deje a su proveedor actual por uno nuevo es el momento de ruptura en la industria del hardware. No se logra con una sola demostración, sino con la acción conjunta del costo, el desempeño, la confiabilidad, la capacidad de entrega y las señales de la industria. Que el cliente cambie de proveedor significa que cree que la nueva solución no solo es más barata, sino suficientemente confiable.

\begin{figure}[H]
\centering
\includegraphics[width=0.96\textwidth]{customer-defection.png}
\caption{Los clientes empiezan a cambiar de proveedor: el momento de ruptura en la cadena de suministro de hardware llega cuando el cliente está dispuesto a dejar la solución anterior por una nueva. Diagrama conceptual propio, elaborado a partir de la conversación de 01:38:15--01:58:07.}
\end{figure}

\begin{knowledgebox}{Por qué los clientes cambian de proveedor}
\begin{tabularx}{\textwidth}{p{0.20\textwidth}p{0.34\textwidth}X}
\toprule
Condición & Qué ve el cliente & Qué significa para el proveedor \\
\midrule
Desempeño confiable & Cumple con los escenarios críticos & La tecnología ya no es solo una demo. \\
Ventaja de costo & El costo total baja de forma notable & Tiene la oportunidad de entrar en modelos de producción en serie. \\
Entrega estable & Puede suministrar según lo programado & Se valida la capacidad de la organización. \\
Riesgo de sustitución controlable & Migrar desde la solución anterior no se sale de control & El cliente está dispuesto a asumir el riesgo del cambio. \\
\bottomrule
\end{tabularx}
\end{knowledgebox}

\subsection{Resumen de la sección}

Para pasar de «se acabó» a que los clientes cambien de proveedor, lo decisivo es pasar de la confianza en la tecnología propia a la validación por parte del cliente. El verdadero punto de inflexión de una empresa de tecnología profunda llega cuando el cliente está dispuesto a confiarte su riesgo.
\section{Hoja de ruta del emprendimiento en tecnología profunda: del fundador técnico a eslabón de la cadena industrial}

Las secciones anteriores trataron por separado la reducción de costos, la estructura accionaria, el financiamiento, los grandes pedidos, la fijación de precios y la deserción de clientes; esta sección integra esos contenidos en una sola hoja de ruta. En la historia de Li Yifan (李一帆), la capacidad técnica es solo el primer paso; la empresa debe incorporar de forma gradual criterio comercial, comprensión del cliente, capacidad de cadena de suministro, ritmo de capitalización y el sistema de las armadoras. El emprendedor en tecnología profunda no pasa del laboratorio a gigante industrial de un solo salto, sino que en cada etapa corrige una carencia.

\begin{knowledgebox}{Etapas del emprendimiento en tecnología profunda}
\begin{tabularx}{\textwidth}{p{0.18\textwidth}p{0.34\textwidth}X}
\toprule
Etapa & Tarea principal & Riesgo típico \\
\midrule
Etapa de laboratorio & Demostrar que la tecnología es viable & Autocomplacencia; no encontrar un mercado real. \\
Etapa de prototipo & Que el cliente vea los resultados & Solo se puede demostrar, no entregar. \\
Etapa de pedidos & Obtener presupuesto real & Los grandes pedidos traen flujo de efectivo y presión de cumplimiento. \\
Etapa de producción en serie & Reducción de costos, rendimiento de fabricación, cadena de suministro y servicio posventa & Los problemas de calidad destruyen la confianza. \\
Etapa de armadoras & Entrar en los modelos de vehículo y en el sistema de suministro a largo plazo & Ciclos de certificación largos y fuerte presión en la negociación. \\
Etapa de globalización & Enfrentar culturas distintas de capital y de clientes & El lenguaje de la confianza y el cumplimiento normativo son complejos. \\
\bottomrule
\end{tabularx}
\end{knowledgebox}

\begin{importantbox}{Lo que el emprendedor en tecnología profunda debe aprender}
El emprendedor de internet suele aprender primero crecimiento y producto; el emprendedor en tecnología profunda suele partir de la tecnología. Lo más difícil de aprender para este último son los clientes, la cadena de suministro, el capital, la producción en serie y la comunicación internacional. Los 11 años de historia de Hesai (禾赛) son precisamente una trayectoria de aprendizaje continuo.
\end{importantbox}

\subsection{Resumen de la sección}

La hoja de ruta del emprendimiento en tecnología profunda muestra que una oportunidad industrial no se materializa por sí sola a partir de una única tecnología. El emprendedor debe llevar la tecnología al prototipo, a los pedidos, a la producción en serie, a las armadoras y a la globalización, y cada paso exige una capacidad nueva.
\section{La entrada al corazón de la industria automotriz: del prototipo al sistema de las armadoras}

La sección anterior trató la deserción de los clientes; esta sección trata la entrada al corazón de la industria automotriz. La cadena de la industria automotriz exige muchísimo a sus proveedores: grado automotriz, confiabilidad, certificación, costo, producción en serie, servicio posventa y suministro a largo plazo deben cumplirse por completo. Pasar del prototipo de LiDAR al sistema de las armadoras no es «vender un componente de hardware», sino convertirse en parte de los sistemas de seguridad e inteligencia del vehículo completo.

\begin{figure}[H]
\centering
\includegraphics[width=0.96\textwidth]{auto-maincamp.png}
\caption{La entrada al corazón de la industria automotriz: para pasar de prototipo técnico a proveedor de grado automotriz hay que entrar en el sistema de las armadoras. Figura conceptual de elaboración propia, a partir de la conversación de 01:58:07--02:38:34.}
\end{figure}

\begin{knowledgebox}{Lectura de la figura: por qué es difícil entrar en el sistema de las armadoras}
Lo que las armadoras buscan es un suministro estable a largo plazo, no un deslumbramiento técnico momentáneo. El proveedor debe demostrar que su producto es confiable, que su costo está bajo control, que sus entregas son estables y que puede sostener el servicio posventa, además de seguir iterando a lo largo del ciclo de vida de cada modelo de vehículo.
\end{knowledgebox}

\subsection{De dónde vienen las oportunidades en la cadena de la industria automotriz}

La sección anterior explicó por qué es difícil entrar en el sistema de las armadoras; esta sección pregunta de dónde vienen, en realidad, las oportunidades de una industria tan difícil. Li Yifan (李一帆) habla una y otra vez de oportunidades. Las oportunidades en la cadena de la industria automotriz no las crea de la nada un emprendedor aislado, sino que surgen de la acción conjunta del auge de los vehículos de nueva energía, la sustitución local de la cadena de suministro, la demanda de inteligencia, la competencia entre armadoras y la capacidad industrial del país. Hesai (禾赛) se encontraba justamente en el punto donde confluyen todas estas oportunidades.

\begin{importantbox}{La estructura de la oportunidad industrial}
\[
\text{Opportunity} \approx \text{Technology} \times \text{Supply Chain} \times \text{Market Timing} \times \text{Customer Adoption}
\]
Aquí, Technology es la madurez técnica, Supply Chain es la viabilidad de la cadena de suministro, Market Timing es la ventana de oportunidad de la industria y Customer Adoption es la disposición de las armadoras a adoptar el producto.
\end{importantbox}

\begin{knowledgebox}{Desglose de las oportunidades en la cadena de la industria automotriz}
\begin{tabularx}{\textwidth}{p{0.20\textwidth}p{0.34\textwidth}X}
\toprule
Origen de la oportunidad & Manifestación & Significado para una empresa de LiDAR \\
\midrule
Crecimiento de los vehículos de nueva energía & Nuevos modelos, nuevas plataformas, nuevas cadenas de suministro & Posibilidad de entrar en el diseño de la siguiente generación de vehículos. \\
Competencia en inteligencia & Aumenta la demanda de asistencia a la conducción, Robotaxi y robots & Se revalúa el valor de los sensores de percepción. \\
Cadena de suministro nacional & Las armadoras están dispuestas a buscar proveedores locales & Abre una ventana de entrada para empresas nuevas. \\
Reducción de costos & Las configuraciones de gama alta llegan a más modelos & Se amplía el tamaño del mercado. \\
Demanda global & Los clientes extranjeros necesitan cadenas de suministro alternativas & Hay que construir confianza internacional y experiencia en grado automotriz. \\
\bottomrule
\end{tabularx}
\end{knowledgebox}

\begin{importantbox}{De empresa de prototipos a proveedor de grado automotriz}
Una empresa de prototipos vende posibilidades; un proveedor de grado automotriz vende certidumbre a largo plazo. La primera necesita deslumbrar con su técnica; el segundo necesita un sistema de calidad, cadena de suministro, servicio posventa, una curva de costos y la confianza de los clientes.
\end{importantbox}

\subsection{Resumen de la sección}

Entrar al corazón de la industria automotriz significa que la empresa deja de ser un proveedor de tecnología y se convierte en un nodo de la cadena industrial. A partir de ese momento, la verdadera competencia ya no es solo de desempeño, sino de grado automotriz, costo, entregas y confianza de las armadoras.
\section{Dinero nuevo, dinero viejo y oportunidades para el país y la nación}

La sección anterior trató la entrada de la empresa en el sistema de las armadoras; esta sección pasa a una narrativa más amplia sobre la globalización y la industria. Li Yifan (李一帆) habla del dinero nuevo y del dinero viejo, y también de las oportunidades para el país, la nación y la industria. El emprendimiento en tecnología dura es distinto del internet de consumo: con frecuencia está ligado a la capacidad manufacturera del país, a la integridad de la cadena industrial, a la comunicación con clientes de todo el mundo y a las diferencias entre culturas de capital.

\begin{figure}[H]
\centering
\includegraphics[width=0.94\textwidth]{new-money-old-money.png}
\caption{Dinero nuevo y dinero viejo: en la comunicación global, cada cultura de capital y de industria tiene su propio lenguaje de confianza. Diagrama conceptual propio, elaborado a partir de la conversación entre 02:38:34 y 03:02:16.}
\end{figure}

\begin{knowledgebox}{Lectura de la figura: globalizarse no es solo traducir el idioma}
El dinero viejo valora la historia, el orden y la estabilidad; el dinero nuevo valora la velocidad, el crecimiento y la tecnología. Cuando una empresa china de tecnología dura se globaliza, no solo tiene que explicar su producto, sino también demostrar que es confiable, que piensa a largo plazo y que puede cumplir sus entregas.
\end{knowledgebox}

\subsection{País, nación y oportunidades industriales}

Después de la globalización, esta subsección vuelve a la narrativa industrial más amplia. Tanto al principio como al final de la conversación aparece la expresión «país, nación, oportunidad». No se trata de una consigna abstracta, sino de una realidad industrial: los vehículos de nueva energía, los sensores, los chips, la cadena de suministro, la manufactura y el mercado global determinan en conjunto si una empresa de tecnología dura tiene una oportunidad. Algunas oportunidades no las crea un emprendedor individual, sino que las otorgan la época y la industria.

\begin{figure}[H]
\centering
\includegraphics[width=0.92\textwidth]{national-industrial-opportunity.png}
\caption{País, nación y oportunidades industriales: el emprendimiento en tecnología dura suele surgir de la suma de la época, la cadena industrial y la capacidad del país. Diagrama conceptual propio, elaborado a partir de la conversación entre 00:00:04 y 00:00:59 y entre 02:38:34 y 03:02:16.}
\end{figure}

\begin{warningbox}{La gran narrativa no sustituye la capacidad de producto}
Las oportunidades para el país y la nación abren una ventana industrial, no una victoria gratuita. La empresa todavía tiene que aprovechar esa ventana con tecnología, costos, calidad, clientes y capacidad organizacional.
\end{warningbox}

\begin{knowledgebox}{El examen final de una empresa de tecnología dura}
\begin{tabularx}{\textwidth}{p{0.22\textwidth}p{0.34\textwidth}X}
\toprule
Examen & Pregunta & Criterio de aprobación \\
\midrule
Examen técnico & ¿El desempeño es realmente superior? & Pruebas objetivas, escenarios de clientes, confiabilidad a largo plazo. \\
Examen de costos & ¿Puede reducir costos de forma sostenida? & BOM, rendimiento de producción, cadena de suministro y escalamiento. \\
Examen de clientes & ¿Las armadoras están dispuestas a adoptarlo? & Certificación, nominación como proveedor, producción en serie y recompra. \\
Examen de capital & ¿Puede atravesar los ciclos? & Financiamiento, flujo de efectivo, salida a bolsa y confianza del mercado. \\
Examen de globalización & ¿Puede cumplir entregas entre culturas distintas? & Clientes en el extranjero, cumplimiento normativo, marca y red de servicio. \\
\bottomrule
\end{tabularx}
\end{knowledgebox}

\begin{knowledgebox}{Lista de verificación para la comunicación global}
\begin{tabularx}{\textwidth}{p{0.22\textwidth}p{0.34\textwidth}X}
\toprule
Dimensión & Qué hay que demostrar & Materiales típicos \\
\midrule
Credibilidad técnica & Que el desempeño y la confiabilidad no son discurso publicitario & Informes de pruebas, casos de clientes, certificación de grado automotriz. \\
Credibilidad comercial & Que la empresa puede suministrar y dar servicio a largo plazo & Finanzas, capacidad de producción, red de posventa. \\
Credibilidad cultural & Que entiende las culturas de distintos clientes y capitales & Equipo local, estilo de comunicación, compromiso a largo plazo. \\
Credibilidad de cumplimiento & Que cumple las normas locales y los requisitos de la cadena de suministro & Revisión de cumplimiento normativo, procesos de datos y de seguridad. \\
\bottomrule
\end{tabularx}
\end{knowledgebox}

\subsection{Resumen de la sección}

El emprendimiento en tecnología dura es una historia de empresa y, al mismo tiempo, una historia de industria y de capacidad nacional. Para entender a Hesai (禾赛), hay que mirar a la vez los detalles técnicos y la ventana que abrió la época.
\section{Síntesis y ampliación}

Esta sección condensa el EP111 en cinco conclusiones. Primera: el LiDAR es el ojo activo de los automóviles y los robots; su valor reside en la medición directa de distancias y en la redundancia de seguridad. Segunda: la reducción de costos del 99.5\% muestra que la industrialización del hardware depende de rediseñar en conjunto el diseño, los chips, la cadena de suministro y la manufactura. Tercera: emprender en tecnología profunda exige atravesar la estructura accionaria, el financiamiento, los pedidos grandes, la fijación de precios, el flujo de efectivo y la entrega. Cuarta: el verdadero punto de inflexión llega cuando los clientes están dispuestos a abandonar la solución anterior y te permiten entrar en el sistema de los fabricantes de automóviles. Quinta: las oportunidades de la tecnología profunda china surgen de la suma de tecnología, cadena industrial, ventana de mercado y capacidad del Estado.

\begin{importantbox}{Ubicar el EP111 en la serie de IA e internet de Zhang Xiaojun (张小珺)}
El EP111 no es una entrevista sobre modelos grandes, pero pertenece a la cadena industrial del automóvil inteligente y la tecnología profunda. Puede leerse junto con el EP120 sobre Physical AI de Xiaopeng (小鹏), el EP132 sobre los robots de Xinghaitu (星海图) y el EP121 sobre los robots de DeepMind, para entender la base de sensores y de cadena de suministro que hay detrás de la IA, la conducción autónoma y los robots.
\end{importantbox}

\begin{knowledgebox}{Por qué este episodio sigue perteneciendo a la línea de IA e internet}
La conducción autónoma y los robots no se reducen a modelos: también necesitan hardware de percepción, una cadena de suministro de grado automotriz, integración con los fabricantes de automóviles y una curva de costos. El LiDAR es una capa de infraestructura previa a que «el modelo entre en el mundo físico»; sin una percepción confiable, a la Physical AI le resulta muy difícil cerrar el ciclo en carreteras y fábricas reales.
\end{knowledgebox}

\begin{knowledgebox}{Relación con otros episodios}
\begin{tabularx}{\textwidth}{p{0.18\textwidth}p{0.34\textwidth}X}
\toprule
Episodio & Tema & Conexión con el EP111 \\
\midrule
EP120 & Physical AI de Xiaopeng y la fábrica de modelos de conducción autónoma & Explica los escenarios de aplicación del LiDAR una vez integrado en el sistema inteligente del vehículo. \\
EP132 & Robots de Xinghaitu y el ciclo cerrado de robot completo, modelo y datos & Muestra cómo los sensores, el robot completo y el acceso a los datos forman en conjunto una barrera competitiva. \\
EP121 & Robots de DeepMind, modelos del mundo y transferencia entre distintos cuerpos robóticos & Aporta el contexto de investigación sobre percepción robótica y aprendizaje de acciones. \\
EP108 & Yu Kai (余凯) relata la historia de la industria de la IA y el automóvil & Junto con Hesai (禾赛), ofrece la perspectiva de los actores poco visibles de la cadena industrial automotriz. \\
\bottomrule
\end{tabularx}
\end{knowledgebox}

\subsection{Takeaways clave}

\begin{enumerate}
\item El valor central del LiDAR es la medición activa de distancias, no simplemente «un sensor más».
\item La capacidad de reducir costos es la capacidad central de un emprendimiento de hardware y determina si la tecnología llega a la producción en masa.
\item Un pedido grande es un evento de credibilidad, no solo un evento de ingresos.
\item Entrar en el sistema de un fabricante de automóviles implica cumplir con el grado automotriz, la calidad, la entrega, el costo y la confianza a largo plazo.
\item El Estado y la cadena industrial ofrecen oportunidades, pero la empresa debe concretarlas con producto y entregas.
\end{enumerate}

\subsection{Preguntas abiertas}

Las conclusiones ya se presentaron; esta sección conserva algunas preguntas que aún requieren observación. Se relacionan directamente con las barreras competitivas de largo plazo de las empresas de LiDAR y con la posición de las empresas chinas de tecnología profunda en la cadena industrial global.

\begin{enumerate}
\item A medida que los modelos de visión sigan mejorando, ¿cómo cambiará el papel del LiDAR en los automóviles de pasajeros y en los robots?
\item ¿De dónde proviene la barrera competitiva de largo plazo de las empresas de LiDAR: de los chips, los algoritmos, la manufactura, la relación con los clientes o el costo a escala?
\item Cuando las empresas chinas de tecnología profunda se internacionalizan, ¿cómo pueden contar bien, al mismo tiempo, la historia tecnológica y la historia de confianza?
\item Entre el desarrollo propio de los fabricantes de automóviles y la especialización de los proveedores, ¿dónde quedará el límite a largo plazo?
\end{enumerate}

\begin{warningbox}{Lo que tienen en común las preguntas abiertas}
Ninguna de estas preguntas puede decidirla por completo una sola empresa. Dependen del avance de los modelos de visión, de la ruta de la conducción autónoma, de la estrategia de compras de los fabricantes de automóviles, de la cadena de suministro internacional y del ciclo de capital. Las empresas de tecnología profunda deben redefinir continuamente su posición conforme cambian estas variables externas.
\end{warningbox}

\subsection{Lecturas adicionales}

\begin{itemize}
\item Si le interesan el automóvil inteligente y la Physical AI, puede compararlo con la entrevista a Liu Xianming (刘先明), de Xiaopeng, en el EP120.
\item Si le interesa la cadena industrial de los robots, puede compararlo con la entrevista a Gao Jiyang (高继扬), de Xinghaitu, en el EP132.
\item Si le interesan los modelos del mundo y la investigación en robótica, puede compararlo con la entrevista a Tan Jie (谭捷), de DeepMind, en el EP121.
\end{itemize}

\end{document}
